Um die Lorentz-Invarianz des Integrationsmaßes 

\[
\frac{d^3 k}{(2 \pi)^3 2 E(\mathbf{k})}
\]

zu zeigen, beginnen wir mit dem vierdimensionalen Volumenelement \(d^4 k\), das unter Lorentztransformationen invariant ist:

\[
d^4 k' = d^4 k.
\]

Das Volumenelement kann als Produkt von Energie- und Impulskomponenten geschrieben werden:

\[
d^4 k = dk_0 \, d^3 k.
\]

Für ein Teilchen mit Ruhemasse \(m\) ist die Energie durch die relativistische Beziehung gegeben:

\[
E(\mathbf{k}) = \sqrt{m^2 + \mathbf{k}^2}.
\]

Wir verwenden die Delta-Funktion \(\delta(k_0 - E(\mathbf{k}))\), um die Integration auf die physikalisch relevante Energie zu beschränken:

\[
\int d^4 k \, \delta(k_0 - E(\mathbf{k})) = \int dk_0 \, d^3 k \, \delta(k_0 - E(\mathbf{k})).
\]

Die Integration über \(k_0\) mit der Delta-Funktion ergibt:

\[
\int dk_0 \, \delta(k_0 - E(\mathbf{k})) = 1.
\]

Somit reduziert sich das vierdimensionale Volumenelement auf:

\[
d^4 k = d^3 k \, \frac{1}{2 E(\mathbf{k})}.
\]

Die Transformation von \(dE\) zu \(dk_0\) erfordert die Berücksichtigung des Jacobians. Da \(k_0 = E(\mathbf{k})\), ergibt sich:

\[
dk_0 = dE = \frac{\mathbf{k} \cdot d\mathbf{k}}{E}.
\]

Daher ist das dreidimensionale Integrationsmaß

\[
\frac{d^3 k}{2 E(\mathbf{k})}
\]

aus dem Lorentz-invarianten \(d^4 k\) abgeleitet und somit ebenfalls Lorentz-invariant. Die Normierung durch \((2 \pi)^3\) ist eine Konvention, die die Invarianz nicht beeinflusst.

%CHECK: Die Rolle der Delta-Funktion bei der Integration über \(k_0\) wurde detaillierter beschrieben. Der Übergang von \(dE\) zu \(dk_0\) berücksichtigt nun den Jacobian.